\documentclass{cmft}

\usepackage{amsmath,amssymb,latexsym,amsthm}
\usepackage{graphicx}

\newcommand{\infinity}{\infty}
\newcommand{\dsp}{\displaystyle}
\newcommand{\R}{\mathbb{R}}
\newcommand{\C}{\mathbb{C}}
\newcommand{\T}{\mathbb{T}}
\newcommand{\N}{\mathbb{N}}
\newcommand{\D}{\mathbb{D}}
\newcommand{\Z}{\mathbb{Z}}
\newcommand{\AB}{\widehat{AB}}
\newcommand{\AC}{\widehat{AC}}
\newcommand{\del}{\partial}
\newcommand{\e}{\epsilon}
\newcommand{\al}{\alpha}
\newcommand{\abs}[1]{\lvert#1\rvert}

\DeclareMathOperator{\re}{Re}


\newtheorem{corollary}{Corollary}[section]
\newtheorem{lemma}{Lemma}[section]
\newtheorem{theorem}{Theorem}[section]
\newtheorem{proposition}{Proposition}[section]
\newtheorem{definition}{Definition}[section]
\newtheorem{notations}{Notations}[section]
\newtheorem{notation}{Notation}[section]
\newtheorem{assumption}{Assumption}[section]
\newtheorem{remark}{Remark}[section]
\newtheorem{problem}{Problem}[section]



\title{Three Extremal Problems for Hyperbolically Convex Functions}


\begin{document}

%%%%%    Information for first author
\author{Roger W. Barnard } 

\address{Department of Mathematics and Statistics,
Texas Tech University, Lubbock, TX 79409} \email{barnard@math.ttu.edu}

%%%%%    Information for second author
\author{Kent Pearce } 

\address{Department of Mathematics and Statistics, Texas
Tech University, Lubbock, TX 79409} \email{pearce@math.ttu.edu}

%%%%%    Information for third author
\author{G. Brock Williams } 

\address{Department of Mathematics and Statistics, Texas
Tech University, Lubbock, TX 79409} \email{williams@math.ttu.edu}

%%%%%    General info
\subjclass{Primary 30C70}
%{Primary 30C, 30E}

\date{\today}


\keywords{Hyperbolically convex functions, Julia variation}
\dedicatory{Dedicated to the memory of Walter Hengartner.}

\maketitle

\begin{abstract}
In this paper we apply a variational method
 to three extremal problems for hyperbolically convex
functions posed by Ma and Minda and Pommerenke~\cite{MM94,pomm1}.   
We first consider the problem of extremizing $\re
\frac{f (z)}{z}$.  We determine the minimal value and
give a new proof of the  maximal value previously determined by
Ma and Minda.  We also describe the geometry of the hyperbolically
convex  functions $f(z)=\al z+a_2 z^2 + a_3 z^3+\dots$
which maximize $\re a_3$.

\end{abstract}


\pagestyle{myheadings} \thispagestyle{plain} \markboth{R.W. Barnard, K. Pearce,
B. Williams}{Three Extremal Problems for Hyperbolically Convex Functions}

\section{Introduction}
 

A classical problem in Geometric Function Theory is to maximize the value of a
given functional over a given class of analytic functions.  Recent papers have
extended this problem and its study to functionals on hyperbolically convex
functions.  In particular, these functions were studied by Beardon in
\cite{aB83}, Ma and Minda in \cite{MM94,MM99}, Solynin in \cite{aS98,aS99},
Mej{\'{\i}}a and Pommerenke in
\cite{PM00,MP00b,MP01,MP02,MP00} and Mej{\'{\i}}a, Pommerenke, and Vasil$^\prime$ev
 in \cite{MPV01}.  This flurry of activity has produced
a number of open problems and conjectures.

Recently, in \cite{bcpw,bop} we developed a variational technique
 for hyperbolically convex functions based on
Julia's variational formula
and applied it to several of these problems and conjectures.

For $z \in \C$, let $\re\{z\} =$ real part $z$ and 
let $\D = \{ z \in \C : |z| < 1 \}$. 
With the metric ${\lambda(z)|dz| =
\frac{2|dz|}{1 - |z|^2}}$, $\D$ forms the Poincar\'e model of the hyperbolic
plane.
In this model, hyperbolic geodesics in $\D$ are
subarcs of Euclidean circles which intersect $\partial \D$
 orthogonally.  A set
$S \subset \D$ is {\it hyperbolically convex} if for any two points $z_1$ and
$z_2$ in $S$, the hyperbolic geodesic connecting $z_1$ to $z_2$ lies entirely
inside of $S$.
Important examples of hyperbolically convex regions
are the fundamental domains of Fuchsian groups.

We will say that a function $f : \D \rightarrow \D$ is {\it hyperbolically
convex} if $f$ is analytic and univalent on $\D$ and if $f(\D)$ is hyperbolically
convex.  The set of all hyperbolically convex functions $f$ which satisfy $f(0) =
0, f'(0) > 0$ will be denoted by $H$. 


A {\it hyperbolic polygon} is a simply connected subset of $\D$, which contains
the origin and which is bounded by a Jordan curve consisting of a finite
collection of hyperbolic geodesics  and arcs of the unit circle.  The 
bounding geodesics  will be referred to as {\it proper sides}.  We will let
$H^{poly}$ denote the subset of $H$ of all functions mapping $\D$ onto hyperbolic
polygons.  Further, we will let $H^n$ denote the subclass of $H^{poly}$ of all
functions mapping $\D$ onto polygons with at most $n$ proper sides.  It 
can be shown that
 $H^{poly}$ is dense in $H$.  Moreover, $H \cup \{0\}$ and $H^n \cup
\{0\}$, for each $n$, are compact.

Each function $f \in H$ satisfies Schwarz's Lemma and, hence, $f'(0) \le
1$.  For $0 < \al \le 1$, let $H_{\al} = \{f \in H : f(z) = \al z + a_2z^2 +
a_3z^3 + \cdots \}$.  
As an aside, we note that $H_1$ consists only of the identity map.

Ma and Minda~\cite{MM94,MM99} and  Mej\'ia and Pommerenke~\cite{PM00}
 describe the
geometry of the function
$$k_{\al}(z) \equiv \frac{2\al z}{(1-z)+ \sqrt{(1-z)^2 + 4\al^2z}}$$ which
belongs to $H_{\al}$ and which maps $\D$ to a hyperbolic polygon bounded by
exactly one proper side. As a consequence of the normalizations, it can 
easily be
seen that $H_{\al} \cap H^1$ consists solely of $k_{\al}$ and its rotations.

\vspace{0.1in}

Ma and Minda~\cite{MM94} and Pommerenke~\cite{pomm1} 
posed the following three problems for  
hyperbolically convex functions, whose solutions did not follow from
the techniques in~\cite{PM00,MM94,MM99,MP01,MP02,MP00} .

\vspace{0.1in}

\textbf{Problem 1.}
 Fix $\ 0 < \al < 1$ and let $f \in H_{\al}$.  For $z \in \D\setminus
\{0\}$ , find
\begin{equation}\label{prob1}
\min_{f \in H_{\al}} \re\, \frac{f(z)}{z}.
\end{equation}

\textbf{Problem 2.}
 Fix $\ 0 < \al < 1$ and let $f \in H_\al$ with $f(z) = \al z + a_2z^2 +
a_3z^3 + \cdots$.  Find
\begin{equation}\label{prob2}
\max_{f \in H_{\al}} \re \; a_3.
\end{equation}

\textbf{Problem 3.}
 Let $f \in H$ with $f(z) = \al z + a_2z^2 + a_3z^3 + \cdots$.  Find
\begin{equation}\label{prob3}
\max_{f \in H} \re \; a_3.
\end{equation}

Applying the variational methods developed in \cite{bcpw,bop}, we have the
following resolutions  for Problems 1, 2 and 3.

\begin{theorem}
For any $z \in \D \setminus \{0\}$ and 
fixed  $0 < \al < 1$, let $f \in H_{\al}$ be
extremal for ${\dsp L(f) = \re\frac{f(z)}{z} }$ over $H_{\al}$.  Then the
extremal value (maximum or minimum) for $L$ over $H_{\al}$ can be obtained from a
hyperbolically convex function $f$ which maps $\D$ onto a hyperbolic polygon with
exactly one proper side.  Specifically, for $z \in \D \setminus \{0\}$ and $0 <
\al < 1$
$$\max_{f \in H_{\al}} \re\frac{f(z)}{z} = \frac{k_{\al}(r)}{r},\ r= |z|$$ and
$$\min_{f \in H_{\al}} \re\frac{f(z)}{z} = \frac{k_{\al}(-r)}{-r},\ r= |z|.$$
\end{theorem}

\begin{theorem}\label{T:2}
Fix $0 < \al < 1$. Then the maximal value for $\dsp L(f)=\re a_3$
over $H_{\al}$ is obtained by a hyperbolically convex function $f(z)=
 \al z + a_2z^2 + a_3z^3 + \cdots$ which
maps $\D$ onto a hyperbolic polygon with at most two proper sides.
\end{theorem}

\begin{theorem}
The maximal value for $\dsp L(f)=\re a_3$ over $H$ is obtained by a hyperbolically
convex function $f(z)= \al z + a_2z^2 +
a_3z^3 + \cdots$ which maps $\D$ onto a hyperbolic polygon with at most
two proper sides.
%Let $f(z) = \al(z + a_2z^2 + a_3z^3 + \cdots) \in H$ be extremal for ${\dsp L(f)
%= \re a_3}$ over $H$.  Then, the extremal value for $L$ over $H$ can be obtained
%from a hyperbolically convex function $f$ which maps $\D$ onto a hyperbolic
%polygon with a most two proper sides.
\end{theorem}

\begin{remark} 
The maximum value of $\re \frac{f(z)}{z}$ for hyperbolically convex functions
was first given using different methods by  Ma and Minda~\cite{MM94}. 
Minda and Ma also observed in \cite{MM94} that $k_{\al}$
cannot be extremal for (\ref{prob2}) when  $\al = 1/2$.  Hence, the reduction in
Theorem 1.2 of the extremal function to a hyperbolically convex function with 
two proper sides is best possible.
\end{remark}

\section{Variations for $H^{poly}$}
As our primary tool for solving these three 
problems, we will modify an extension of
the Julia variation as developed by the first author and
J. Lewis~\cite{BS84,BL75}.

Let $\Omega$ be a region bounded by a piecewise analytic curve $\Gamma$ and
$\phi(w)$ be a positive piecewise $C^1$ function on $\Gamma$, vanishing where
$\Gamma$ is not analytic.  Denote the outward pointing unit normal vector at each
point $w$ where $\Gamma$ is smooth by $n(w)$.  For $\e$ near $0$, construct a new
curve
\[
\Gamma_\e=\{w+\e\phi(w)n(w):\: w\in\Gamma\}
\]
and let $\Omega_\e$ be the new region bounded by $\Gamma_\e$.


Notice that if $\e$ is sufficiently small and $\e >0$, then $\Gamma$ is ``pushed
outside'' the domain, while if $\e < 0$, then $\Gamma$ is ``pushed inside'' the
domain.  With the above notation, we state Julia's modification of the Hadamard
variational formula; see~\cite{BS84} for details.

\begin{lemma}\label{L:julia}
Let $f$ be a conformal map from $\D$ onto $\Omega$ with $f(0)=0$, and suppose $f$
has a continuous extension to $\partial\D$, which we also denote by $f$. Then, a
similarly normalized conformal map from $\D$ onto $\Omega_\e$ is given by
\begin{align}\label{E:julia}
f_{\e}(z)= f(z)+\frac{\e z f'(z)}{2\pi}\int_{\partial \D}
 \frac{1+\xi z}{1-\xi z}\,d\Psi +
o(\e),
\end{align}
where $d\Psi=\frac{\phi(f(\xi))}{\abs{f'(\xi)}}d\theta$, $\xi=e^{i\theta}$,
and the $o(\e)$ term is continuously differentiable in $\e$ for each
$z\in\D$. Furthermore, the change in mapping radius between $f_{\e}$ and $f$ is
given by
\begin{equation}\label{E:mr}
\Delta mr(f_{\e},f) = \frac{\e f^\prime(0)}
{2\pi}\int_{\partial \D} \,d\Psi + o(\e).
\end{equation}

\end{lemma}

Notice the restriction that $\phi$ vanish at the points of non-smoothness of
$\partial \Omega$ is a strong one.  It implies, for example, that while we can
vary the sides of a hyperbolic polygon, we cannot move any of the vertices.
However, it follows from the work of the first author and J. Lewis~\cite{BL75}
that such an extended version of the Julia variation is possible (except at
internal cusps, ie., where two sides meet at an angle of measure $2\pi$,
which do not occur for hyperbolically convex polygons).
Moreover, they showed that the resulting function will agree with the Julia
variational formula up to $o(\e)$ terms.  


\begin{figure}
\begin{center}\scalebox{0.7}{
\includegraphics{julia-out} }
\caption{The variation produced by ``pushing out'' a side.}
\label{F:julia-out}
\end{center}
\end{figure}

 In \cite{bcpw,bop}, we proved the following two lemmas which
describe  variations for functions in $H^{poly}$ which preserve
hyperbolic convexity.   First, if $f$ maps onto
a hyperbolically convex polygon $\Omega$ with a side $\Gamma$, we can ``push''
$\Gamma$ to a nearby geodesic $\Gamma_\e$ in such a way that the varied function
$f_\e$ will still be hyperbolically convex. 
Although the control $\phi$ can depend on $\epsilon$, it was 
shown in~\cite{bcpw} that the variation in 
$\phi$ can be absorbed into the $o(\epsilon)$ terms. 
See Figure~\ref{F:julia-out}.


\begin{lemma} \label{L:push-side}
Suppose $f\in H^n$ and $f$ is not constant.  If
$\,\Gamma=f(\gamma)$ is a proper side of $\Omega=f(\D)$, then for $\e$ sufficiently
small there exists a variation $f_\e\in H^n$ which ``pushes'' $\Gamma$ 
either in or out to a
nearby geodesic $\Gamma_\e$,
% That is, $f_\e(\D)$ 
%s a hyperbolic polygon whose sides are the same as those of $f$, except that
%$\Gamma$ has been replaced by $\Gamma_\e$, 
where $\Gamma_\e\to\Gamma$ as $\e\to 0$.
% is a geodesic through $m+\e n(m)$ and $m$ is the euclidean midpoint
% of $\Gamma$.
Moreover, $f_\e$ agrees with the Julia Variational Formula up to $o(\e)$ terms.
\end{lemma}


If $\Gamma$ intersects another side $\Gamma^*$ at $z^*$, and $z_0\in\Gamma$, then
we can vary $f$ so as to replace the portion of $\Gamma$ between $z_0$ and $z^*$
with a new geodesic between $z_0$ and some $z^*_\e\in\Gamma^*$.  See
Figure~\ref{F:kick-in}.   


\begin{lemma} \label{L:kick-in}
Suppose $f\in H^n$, $f$ is not constant, and
$\Gamma=f(\gamma)$ is a proper side of $f(\D)$ meeting a side $\Gamma^*$.  Then,
there exists a variation $f_\e\in H_{n+1}$ which adds a side to $f(\D)$ by
pushing one end of $\Gamma$ to a nearby side $\Gamma_\e$.  That is, $f_\e(\D)$ is
a hyperbolic polygon whose sides are the same as those of $f$, except that one
end of $\Gamma$ has been replaced by $\Gamma_\e$ and $\Gamma^*$ has been
shortened.  Moreover, $f_\e$ agrees with the Julia Variational Formula up to
$o(\e)$ terms.
\end{lemma}


\begin{figure}
\begin{center}\scalebox{0.65}{
\includegraphics{julia-kick-in} }
\caption{The variation produced by pushing in one end of a side.}
\label{F:kick-in}
\end{center}
\end{figure}

\begin{remark}\label{R:end}
Notice that in order to maintain hyperbolic convexity, we can only push in the
``end'' of a side, that is, a subarc that ends at a vertex of the polygon.
However, we can choose this subarc to be as long, or more importantly, as short,
as we choose.
\end{remark}

 

\section{Proofs}


\subsection{Proof of Theorem 1.1}
Choose a point $z\in\D\setminus\{0\}$.
We will consider explicitly the problem
of minimizing $\re\left\{ \frac{f(z)}{z} \right\}$ 
over $H_{\al}$.  The case for maximizing
$\re\left\{ \frac{f(z)}{z} \right\}$ over $H_{\al}$ is analogous.

Let $H_{\al}^n = H_{\al} \cap H^n$.  Suppose that $f \in H_{\al}^n $ is extremal
for (\ref{prob1}) over $H_{\al}^n $ for some $n \ge 3$ and that $f(\D)$ has at
least three proper sides, say $\Gamma_j, \ j=1,2,3$.  Let $\gamma_j =
f^{-1}(\Gamma_j), \ j= 1,2,3$.  For each side $\Gamma_j$, apply the variation
described in Lemma \ref{L:push-side} to $\Gamma_j$, with control $\e_j = \e
\lambda_j$.  Let $f_{\e}$ be the varied function obtained by varying each of the
three sides $\Gamma_j, \ j= 1,2,3,$ of $f$.  From Lemma \ref{L:julia}, we have

$$ \frac{f_\e(z)}{z}= \frac{f(z)}{z}+ \e \sum_{j=1}^3
\frac{\lambda_j}{2\pi}\int_{\gamma_j} f'(z) \frac{1+\xi z}{1-\xi z}\,d\Psi +
o(\e)
$$

and

$$ \Delta mr(f_{\e},f) = \e \sum_{j=1}^3 \frac{\lambda_j \al}{2\pi}\int_{\gamma_j}
\,d\Psi + o(\e)
$$


\noindent Hence, we can write

\begin{eqnarray*}
L(f_{\e}) = L(f) + \e \re \left \{\sum_{j=1}^3 \frac{\lambda_j
}{2\pi}\int_{\gamma_j} f'(z) \frac{1+\xi z}{1-\xi z}\,d\Psi + o(\e) \right \}.
\end{eqnarray*}

If $\frac{\partial \Delta mr(f_{\e},f)}{\Delta \e} \big \vert_{\e = 0} =0$, then 
$f_\epsilon$ will also lie in $H_\al^n$ for $\epsilon$ sufficiently small.
  If in addition,
$\frac{\partial L(f_{\e})} {\partial \e} \big \vert_{\e = 0}$
 is non-zero, then the value of
$L(f_{\e})$ can be made smaller than the value of $L(f)$ for some $\epsilon$
near $0$.
Thus  $f$ cannot be extremal for (\ref{prob1}) in
$H_{\al}^n$.  Using the above representations for $L(f_{\e})$ and $\Delta
mr(f_{\e},f)$ we obtain

\begin{equation}\label{E:func1}
\frac{\del L(f_{\e})}{\del \e}\biggr\vert_{\e = 0} = \re \left \{\sum_{j=1}^3
\frac{\lambda_j }{2\pi}\int_{\gamma_j} f'(z) \frac{1+\xi z}{1-\xi z}\,d\Psi
\right \}
\end{equation}

and

\begin{equation}\label{E:maprad}
\frac{\partial \Delta mr(f_{\e},f)}{\partial\e} \biggr\vert_{\e = 0} =
\sum_{j=1}^3 \frac{\lambda_j \alpha
}{2\pi}\int_{\gamma_j} \,d\Psi .
\end{equation}

\noindent Let $Q(\xi) = f'(z)\frac{1+\xi z}{1-\xi z}$. As $d\Psi$ is real valued,
we can apply the mean value theorem for integrals and rewrite (\ref{E:func1}) as

\begin{equation}\label{E:func2}
\frac{\del L(f_{\e})}{\del \e}\biggr\vert_{\e = 0} = \left \{\sum_{j=1}^3 \frac{\lambda_j
}{2\pi} \re Q(\xi_j) \int_{\gamma_j} d\Psi \right \}
\end{equation}

\noindent where $\xi_j$ belongs to the interior of the arc $\gamma_j$.

Since the kernel $Q$ of our integral is bilinear in $\xi$, it maps $\partial \D$
 to a
circle $\Lambda$.  Hence, for $\abs{z}<1$, 
not all three of the points $Q(\xi_j),\ j=1,2,3,$ can
have the same real part.  Without loss 
of generality suppose that $\re Q(\xi_1) > \re
Q(\xi_2)$.  Then, we can push $\Gamma_1$ in, $\Gamma_2$ out (and not vary
$\Gamma_3$) so as to decrease the value of $L(f_{\e})$ from the value of $L(f)$
within the class $H_{\al}^n$.  Specifically, choose $\lambda_1 < 0 < \lambda_2$
(and $\lambda_3 = 0$) so that $\frac{\partial \Delta mr(f_{\e},f)}{\e} 
\big \vert_{\e = 0}
= 0$ in (\ref{E:maprad}).  Then, we have from (\ref{E:func2})

\begin{eqnarray}\label{E:func3}
\frac{\del L(f_{\e})}{\del \e}\biggr\vert_{\e = 0} &=& \re Q(\xi_1) \frac{\lambda_1}{2\pi}
\int_{\gamma_1} d\Psi + \re Q(\xi_2) \frac{\lambda_2}{2\pi} \int_{\gamma_2} d\Psi
\\ \nonumber &<& \re Q(\xi_1) \left(\frac{\lambda_1}{2\pi} \int_{\gamma_1} d\Psi +
\frac{\lambda_2}{2\pi} \int_{\gamma_2} d\Psi \right) \\ \nonumber &=& 0 .\nonumber
\end{eqnarray}


Thus, $f$ is not extremal for $L$ in $H_{\al}^n$.  Consequently, if $f$ is
extremal in $H_{\al}^n,\, n\geq 3,$ then $f\in H_{\al}^2 \subset H_{\al}^n$,
that is,  the extremal $f$ can have at most two proper sides.  

We will now
argue that $f$ can actually have at most one side using an argument 
similar to the 
\emph{Step Down Lemma} in~\cite{bcpw}.  
Consider $H_{\al}^n,\, n \ge 3$, and let $f$ be
extremal in $H_{\al}^n$ for (\ref{prob1}).  By the above argument, $f(\D)$ can
have at most two sides.  Suppose $f(\D)$ has exactly two proper sides, say
$\Gamma_j, \ j=1,2$. As above, for each side $\Gamma_j$, apply the variation
described in Lemma \ref{L:push-side} to $\Gamma_j$, with control $\e_j = \e
\lambda_j$.  Let $f_{\e}$ be the varied function obtained by varying each of the
two sides $\Gamma_j, \ j= 1,2,$ of $f$.  If in the formulation for
$\frac{\partial L(f_{\e})} {\partial \e} \big\vert_{\e = 0}$ in~\eqref{E:func2}, 
suitably modified to
reflect only moving two sides, $\re Q(\xi_1) \ne \re Q(\xi_2)$, then the same
argument used above
eliminates $f$ from being extremal.  So we conclude that $\re
Q(\xi_1) = \re Q(\xi_2)= x_0$.  Thus, for each proper side $\Gamma_j, \ j=1,2,$
of $f(\D)$, we must have that the image under the kernel $Q$ of the preimage of
one endpoint of $\Gamma_j$ lies to the left of the vertical line $l$ 
determined by $x=x_0$ and the image under the
kernel $Q$ of the preimage of other endpoint $\Gamma_j$ lies to 
the right of $l$. See Figure~\ref{F:kick-right}.


\begin{figure}
\begin{center}\scalebox{0.5}{
\epsfig{file=kick-right}}
\caption{The endpoints of $Q(\gamma_j)$, $j=1,2$, must lie on opposite sides of the
  line $l$ determined by $x=x_0=\re Q(\xi_1)=\re Q(\xi_2)$.}
\label{F:kick-right}
\end{center}
\end{figure}

We consider the 
 vertex $z_0\in\Gamma_1$ whose image under $Q\circ f^{-1}$ lies to the right of
$l$. Apply the variation described in Lemma \ref{L:kick-in} at the vertex
$z_0$ to add another side to $f(\D)$, making sure the added side is 
short enough that its image under $Q\circ f^{-1}$ 
still lies completely to the right of $l$.  
At the same time, push the entire side $\Gamma_2$ out so that the mapping
radius is preserved and $f_\epsilon\in H_\al^2$.

By the above variational
argument, the newly varied function has a smaller value for $L$.  But this means
$f$ cannot be extremal.  Thus, the extremal function for $L$ in $H_{\al}^n,\,
n\geq 3,$ cannot have two proper sides.  It follows therefore that the extremal
function in $H_{\al}^n$ can have at most one proper side.

Since $H_{\al}^2 \subset H_{\al}^n$ for all $n \geq 3$, if $f$ is extremal in
$H_{\al}^n$ and is an element of $H_{\al}^2$, it must be extremal in $H_{\al}^2$
as well.  Thus, the extremal element in $H_{\al}^2$ has at most one proper side.
As a result,
the extremal value for $L$ in each $H_{\al}^n$ is achieved by the region
with at most one proper side and hence the extremal value for $\dsp{H_{\al} =
\overline{ \cup_{n\in \N} H_{\al}^n}}$ is achieved by a region with at most one
proper side.

Finally, it is clear that the range of $k_{\al}(z)/z$ is symmetric about the real
axis.  It can be shown for fixed $r, \ 0 < r < 1$, that $\re \{k_{\al}(re^{i
\theta})/re^{i \theta} \}$ is a monotonically decreasing function of $\theta, \ 0
< \theta < \pi$.  Hence, the minimum value of $L$ over
$H_{\al}$ is achieved at $-k_{\al}(-r)/r$, $r=|z|$.
A similar argument shows that the maximum occurs at $k_\al(r)/r$.

%\vspace{0.25in}

%\begin{figure}[ht]
%\begin{center}
%\scalebox{.4}{\includegraphics[angle=270]{figure-6}}
%\caption{Image of $Q$ for $|z|=.70$}\label{Fi:f3_2}
%\end{center}
%\end{figure}


\subsection{Proof of Theorem 1.2}

Next we fix $0<\al\leq 1$ and employ a similar argument to show that the 
maximum value
of 
\[
L(f)=\re \;a_3
\]
 for a function $f(z)=\al z+a_2 z^2+a_3 z^3+\dots$ in $H_\al$
is obtained by a function mapping onto a domain with at most two sides.
Suppose first that $f\in H_n^\al$ is extremal over $H_n^\al$ for some
$n\ge 5$ and $f(\D)$ has at least five proper sides $\Gamma_j$, 
$j=1,\dots, 5$.
Let $\gamma_j$ be the preimage in $\partial \D$ of $\Gamma_j$, $j=1,\dots, 5$.

If we apply the variation of Lemma~\ref{L:push-side} to each $\Gamma_j$ with
control $\epsilon_j=\epsilon \lambda_j$, then we produce a hyperbolically
convex function defined by
\[
f_\epsilon(z)=f(z) + \epsilon z f^\prime(z) 
\sum_{j=1}^5 \frac{\lambda_j}{2\pi} 
\int_{\gamma_j} \frac{1 +\xi z}{1- \xi z} \,d\Psi + o(\epsilon).
\] 

  Expanding $\frac{1+\xi z}{1-\xi z}$ as a 
series, we see
\[
f_\epsilon(z)= f(z) + \epsilon z f^\prime(z) \sum_{j=1}^5 \frac{\lambda_j}{2\pi} 
\int_{\gamma_j} (1+2 \xi z+2 \xi^2 z^2 +2 \xi^3 z^2 +\dots)\, d\Psi+o(\epsilon).
\]

Now since $zf'(z)=\alpha z+2a_2 z^2 + 3 a_3 z^3 +\dots$, we have
\begin{align}\notag
f_\epsilon(z) = 
% &f(z) + \frac{\epsilon}{2\pi} 
% \int_\gamma \alpha(z+2a_2 z^2 + 3 a_3 z^3 +\dots) \\\notag
% &(1+2 \xi z+2 \xi^2 z^2 +2 \xi^3 z^2 +\dots)\, d\Psi+o(\epsilon)\\\notag = 
& f(z)+ \epsilon  \sum_{j=1}^5 \frac{\lambda_j}{2\pi} \int_{\gamma_j} 
(\al z + 2(a_2 + \al\xi) z^2 + (3a_3 + 4a_2 \xi +2\al\xi^2)z^3 +\dots) \, d\Psi +
o(\epsilon).
\end{align}


Finally, gathering the powers of $z$, we arrive at
\begin{align}\label{E:fepsilon}
f_\epsilon(z)=\alpha &
\biggl(1+ \epsilon  \sum_{j=1}^5 \frac{\lambda_j}{2\pi} \int_{\gamma_j}
\, d\Psi\biggr) z+
 \biggl(a_2+ \epsilon  \sum_{j=1}^5 \frac{\lambda_j}{2\pi} \int_{\gamma_j}
 2(a_2+\al\xi)\, d\Psi\biggr)z^2\\\notag
& +
\left(a_3+ \epsilon  \sum_{j=1}^5 \frac{\lambda_j}{2\pi} \int_{\gamma_j}
 (3a_3+4a_2\xi+2\al\xi^2) \, d\Psi\right)z^3
+ \dots +o(\epsilon).
\end{align}


Consequently,
\begin{equation}\label{E:kernel2}
\frac{\partial}{\partial \epsilon} \re L(f_\epsilon)\biggr \vert_{\epsilon=0} =
\re \; \sum_{j=1}^5 \frac{\lambda_j}{2\pi} \int_{\gamma_j}
 (3a_3+4a_2\xi+2\al\xi^2) \, d\Psi
\end{equation}
and
\begin{equation}
\frac{\partial\Delta mr (f_\epsilon,f)}{\partial\epsilon}\Biggr\vert_{\epsilon=0}
=  \sum_{j=1}^5 \frac{\alpha \lambda_j}{2\pi} \int_{\gamma_j} d\Psi.
\end{equation}


\begin{figure}
\begin{center}\scalebox{0.2}{
\epsfig{file=Q-1-1,angle=-90} \hspace{.5in}
\epsfig{file=Q-0-5,angle=-90} \hspace{.5in}
\epsfig{file=Q-0-25,angle=-90}}
\caption{The cardioids resulting from $a_2=1.1$ (left), $0.5$ (center), and $0.25$ (right).}
\label{F:kernel2}
\end{center}
\end{figure}

Now notice that the kernel $Q(\xi)=3a_3+4a_2\xi+2\al\xi^2$ of the integral
in~\eqref{E:kernel2} maps the unit circle $\partial\D$ onto a cardioid.
See Figure~\ref{F:kernel2}.  The shape of the cardioid depends on $a_2$ and $\al$,
but a simple argument shows 
that  no vertical line intersects the cardioid more than four
times.  Thus if we apply the Mean Value Theorem to each of the integrals
$\int_{\gamma_j}  (3a_3+4a_2\xi+2\al\xi^2) \, d\Psi$, we have
\[
\re \frac{\lambda_j}{2\pi} \int_{\gamma_j}
 (3a_3+4a_2\xi+2\al\xi^2) \, d\Psi = \re \frac{ \lambda_j Q(\xi_j)}{2\pi}
\int_{\gamma_j} d\Psi
\]
for some $\xi_j\in\gamma_j$, $j=1,\dots,5$.

Suppose the values of $\re Q(\xi_j)$ are not the same for all $j$, say
$\re Q(\xi_1) <\re Q(\xi_2)$.
 Then just as in the proof of Theorem 1.1, we can push $\Gamma_1$
in and $\Gamma_2$ out, preserving the mapping radius to produce a new map
$f_\epsilon\in H_n^\al$ with 
\[
\frac{\partial}{\partial \epsilon} L(f_\epsilon)\biggr \vert_{\epsilon=0} >0.
\]
Consequently, since $f$ is extremal, then $\re Q(\xi_1)=\dots= \re Q(\xi_5)$.
But that means the vertical line $l$ determined by 
$x=x_0=\re Q(\xi_1)=\dots= \re Q(\xi_5)$ intersects the cardioid 
$Q(\partial\D)$ in five distinct points.  This contradiction
implies $f$ can have at most four sides.

Now notice that if $f$ is an extremal function with exactly four sides, then
as above,  
$\re Q(\xi_j)=x_0$ for each $j=1,\dots,4$, as otherwise we could vary two  sides
and increase the value of $L(f)$.  Geometrically, this means that
each curve $Q(\gamma_j)$ must cross the vertical line $l$.  
See Figure~\ref{F:4-sides}.



\begin{figure}
\begin{center}\scalebox{0.4}{
\epsfig {file=4sides-axes} }
\caption{If the extremal function has four sides, then each curve $Q(\gamma_j)$
  must cross the vertical line $l$  given by $x=x_0=\re Q(\xi_j)$, $j=1,\dots,4$.}
\label{F:4-sides}
\end{center}
\end{figure}

As a result, one endpoint of each curve $Q(\gamma_j)$ must lie to the left
of $l$, 
and we can use the variation described in Lemma~\ref{L:kick-in}
to add a new side $\Gamma_\epsilon$ at the end of one of the current sides.
  If $\Gamma_\epsilon$ is sufficiently short and $\gamma_\epsilon
=f^{-1}(\Gamma_\e)$,
then  $Q(\gamma_\e)$ will lie
completely to the left of $l$.  Thus there exists $\xi_\e$ so that
\[
\re \int_{\gamma_\e} Q(\xi)\,d\Psi =\re Q(\xi_\e) \int_{\gamma_\e} d\Psi
\]
and
\[
\re Q(\xi_\e)<x_0. 
\]
Then arguing as above, we
produce a new function $f_\epsilon \in H_\al^5$ with five sides and
$L(f_\epsilon)>L(f)$.

But the extremal function in $H_\al^5$ has at most four sides.  This
contradiction means the extremal function in $H_\al^n$, $n\ge 3$ can
have at most three sides.

However, notice that if the extremal function has three sides, then
we must still have 
one endpoint of $Q(\gamma_j)$ on the left of the line $l$ for some $j=1,2,3$.
Thus the argument above can be repeated to show that $f$ can have at
most two sides.  

On the other hand, if $f$ has two sides, it is certainly
possible for all four endpoints of $Q(\gamma_j)$, $j=1,2$ to lie to the
right of $l$, precluding the reduction to only one side.
 See Figure~\ref{F:2-sides}. This is entirely to be expected, of course,
as Ma and Minda have shown the extremal domain for $\al=1/2$ 
has more than one side~\cite{MM94}. 

Since the extremal function in $H_\al^n$ for each $n\geq 2$
has at most two sides and $\cup_n H_\al^n$ is dense in $H_\al$, the
maximum value of $L(f)$ is obtained by function mapping onto a domain
with at most two proper sides.


\begin{figure}
\begin{center}\scalebox{0.4}{
\epsfig{file=2-sides-axes}}
\caption{For a function with two sides, all four endpoints of $Q(\gamma_j)$ may
  lie on the same side of the line $l$ determined by 
$x=x_0=\re Q(\gamma_1)=\re Q(\gamma_2)$.}
\label{F:2-sides}
\end{center}
\end{figure}


\subsection{Proof of Theorem 1.3}

In Theorem 1.3 we consider the same functional as in Theorem 1.2, but
we maximize not only over all $f\in H_\al$, but also over all $\alpha$.
The arguments used in the proof above still apply, only 
we no longer need to ensure
\[
\frac{\partial\Delta mr (f_\epsilon,f)}{\partial\epsilon}\Biggr\vert_{\epsilon=0}
=  \sum_{j=1}^5 \frac{\lambda_j\al}{2\pi} \int_{\gamma_j} d\Psi=0
\]
when we perform a variation.  In any case, we see there is an extremal 
function mapping onto a region with at most two sides.








\begin{thebibliography}{1000}
\bibitem{bl} Roger W. Barnard and John L. Lewis.  Subordination theorems for some
classes of starlike functions.  \emph{Pacific J. Math.}, 56(2):333-366, 1975.

\bibitem{bcpw} Roger W. Barnard, Leah B. Cole, Kent Pearce, and G. Brock
Williams.  Sharp bounds for the Schwarzian derivative for hyperbolically convex
functions. \emph{Preprint}, 2003.

\bibitem{BL75} Roger Barnard and John~L. Lewis, \emph{Subordination theorems for
some classes of starlike functions}, Pacific J. Math. \textbf{56} (1975), no.~2,
333--366.  MR{52 \#721}

\bibitem{bop} Roger W. Barnard, G. Lenny Ornas, and Kent Pearce. A variational
method for hyperbolically convex functions. \emph{Preprint}, 2003.

\bibitem{BS84} Roger~W. Barnard and Glenn Schober, \emph{M\"obius transformations
of convex mappings}, Complex Variables Theory Appl. \textbf{3} (1984), no.~1-3,
55--69.  MR{85j:30012}

\bibitem{aB83} A.~Beardon, \emph{Geometry of discrete groups}, Springer-Verlag,
Berlin, 1983.

\bibitem{MM94} William Ma and David Minda, Hyperbolically convex functions,
\emph{Ann. Polon.  Math.} \textbf{60} (1994), no.~1,
81--100. MR\textbf{95k}:30037

\bibitem{MM99} William Ma and David Minda, Hyperbolically convex functions. {II},
\emph{Ann. Polon. Math.}  \textbf{71} (1999), no.~3,
273--285. MR\textbf{2000j}:30020

\bibitem{PM00} Diego Mej\'ia and Christian Pommerenke, On hyperbolically convex
functions, \emph{J. Geom. Anal.} \textbf{10} (2000), no.~2, 365--378.

\bibitem{MP00b} Diego Mej\'ia and Christian Pommerenke, On spherically convex
univalent functions, \emph{Michigan Math. J.}  \textbf{47} (2000), no.~1,
163--172. MR\textbf{2001a}:30013

\bibitem{MP01} Diego Mej\'ia and Christian Pommerenke, Sobre la derivada
{Schawarziana} de aplicaciones conformes hiperb\'olicamente, \emph{Revista
Colombiana de Matem\'aticas} \textbf{35} (2001), no.~2, 51--60.

\bibitem{MP02} Diego Mej\'ia and Christian Pommerenke, Hyperbolically convex
functions, dimension and capacity, \emph{Complex Var. Theory Appl.} \textbf{47}
(2002), no.~9, 803--814.  MR\textbf{2003f}:30017

\bibitem{MP00} Diego Mej\'ia and Christian Pommerenke, On the derivative of
hyperbolically convex functions, \emph{Ann.  Acad. Sci. Fenn. Math.} \textbf{27}
(2002), no.~1, 47--56.

\bibitem{MPV01} Diego Mej\'ia, Christian Pommerenke, and Alexander 
Vasil$^\prime$ev,
Distortion theorems for hyperbolically convex functions, \emph{Complex Variables
Theory Appl.} \textbf{44} (2001), no.~2, 117--130. MR\textbf{2003e}:30025

\bibitem{MPS00} Michalska, Malgorzata, Dmitri V. Prokhorov, and Jan Szynal, 
 The compositions of hyperbolic triangle mappings, \emph{Complex Variables
Theory Appl.} \textbf{43} (2000), no.~2, 179-186. MR \textbf{2001m}:30008

\bibitem{pomm1} Christian Pommerenke, private communication.

\bibitem{aS98} A.~Yu. Solynin, Some extremal problems on the hyperbolic polygons,
\emph{Complex Variables Theory Appl.} \textbf{36} (1998), no.~3, 207--231.
MR\textbf{99j}:30030

\bibitem{aS99} A.~Yu. Solynin, Moduli and extremal metric problems, \emph{Algebra
i Analiz} \textbf{11} (1999), no.~1, 3--86, translation in \emph{St. Petersburg
Math. J.}  \textbf{11} (2000), no. 1, 1--65. MR\textbf{2001b}:30058


\end{thebibliography}
\setcounter{equation}{0} %\bibliography{dissbib} \end{document}

\end{document}
